% =============================================================================
% layout.tex —— 版式层：页标题、边距、页脚与页码、过渡页、目录与列表模板
% -----------------------------------------------------------------------------
% 依赖：colors.tex（字体颜色）、style-*.tex（\InkMargin、\InkTopRule 等）。
% 本文件是「改版面观感」的主要入口。
% =============================================================================

% -----------------------------------------------------------------------------
% 1. 页标题与边距
% -----------------------------------------------------------------------------
% 标题、横线和正文共享同一组左右边距，只需在 style-*.tex 里改 \InkMargin。
\setbeamersize{text margin left=\InkMargin,text margin right=\InkMargin}

% 标题下方蓝色横线的粗细
\newcommand{\InkTopRule}{0.7pt}

% 自定义页标题模板：标题（+可选副标题）→ 蓝色细横线 → 下方留白。
% 用 \hrule 而不是 \hline，避免额外的表格上下文；
% 顶部 0.35cm 与底部 0.25cm 的留白决定正文起始位置。
\setbeamertemplate{frametitle}{%
  \nointerlineskip
  \begin{beamercolorbox}[wd=\textwidth,sep=0pt]{frametitle}
    \vskip0.35cm
    \usebeamerfont{frametitle}\strut\insertframetitle\strut\par
    \ifx\insertframesubtitle\empty\else
      {\usebeamerfont{framesubtitle}\insertframesubtitle\par}%
    \fi
    \vskip0.12cm
    {\color{AcademicBlue}\hrule height\InkTopRule}
    \vskip0.25cm
  \end{beamercolorbox}}

% 多页 frame 自动续页时，在标题后追加序号，例如「（2）」
\setbeamertemplate{frametitle continuation}{%
  \ifnum\insertcontinuationcount>1\relax\space(\insertcontinuationcount)\fi}

% 标题、章节页用衬线粗体（中文即思源宋体 Bold）
\setbeamerfont{title}{family=\rmfamily,series=\bfseries}
\setbeamerfont{frametitle}{family=\rmfamily,series=\bfseries}
\setbeamerfont{section title}{family=\rmfamily,series=\bfseries}

% 关闭右下角导航图标；文献表不显示条目图标
\setbeamertemplate{navigation symbols}{}
\setbeamertemplate{bibliography item}{}
% 图表编号
\setbeamertemplate{caption}[numbered]

% -----------------------------------------------------------------------------
% 2. 列表模板
% -----------------------------------------------------------------------------
\setbeamertemplate{enumerate item}{\insertenumlabel.}
\setbeamertemplate{enumerate subitem}{\insertenumlabel.\insertsubenumlabel.}
\setbeamertemplate{itemize item}{\small\textbullet}
\setbeamertemplate{itemize subitem}{\textendash}

% -----------------------------------------------------------------------------
% 3. 页脚与页码（含附录独立编号）
% -----------------------------------------------------------------------------
% 正文页码：「当前 / 总」。沿用 Beamer 内置的 totalframenumber 模板，
% 其内容等价于原来的 \insertframenumber\,/\,\inserttotalframenumber。
%
% 附录页码：独立编号「A-序号 / A-总数」。
%   \insertframenumberinappendix   附录内第几页（从 1 开始）
%   \insertappendixframenumber     附录总页数
% 这两个计数器由 Beamer 在 \appendix 时维护，并借助 .nav 辅助文件跨轮传递，
% 因此不需要手工硬编码任何数字，也不需要自己加计数器。
% 注意：绝不能在页脚里判断 \ifbeamer@inappendix —— 该条件在导言区
%       渲染页脚时还未就绪，会直接报 Undefined control sequence。
%       正确做法是用下面的 \appendix 包装，在进入附录时切换模板。
\defbeamertemplate*{page number in head/foot}{inkappendixframenumber}{%
  A-\insertframenumberinappendix\,/\,A-\insertappendixframenumber}

% 当前使用的页码模板名；\appendix 会把它切换成附录模板。
% 使用 \gdef 是因为 \appendix 在正文中执行，需要全局生效。
\def\InkPageNumberStyle{totalframenumber}

% 统一的页码输出：右对齐 + 当前模板。字号/颜色由 \ApplySlideStyle 与 colors.tex 设定。
% \setbeamertemplate 在这里位于每个页脚的分组内部，作用范围仅限本次排版；
% 由于模板名是全局的，切换一次即对所有后续页面稳定生效。
\newcommand{\InkPageNumber}{%
  \setbeamertemplate{page number in head/foot}[\InkPageNumberStyle]%
  \usebeamerfont{page number in head/foot}%
  \usebeamercolor[fg]{page number in head/foot}\hfill
  \usebeamertemplate***{page number in head/foot}\par}

% 包装 \appendix：先执行 Beamer 原本的附录处理（进入附录模式、写入计数器），
% 再把页码模板切换到附录版本。这样附录页脚显示 A-1 / A-3 式编号，
% 而正文始终保持「当前 / 总」不变。
\let\InkAppendixOriginal\appendix
\renewcommand{\appendix}{%
  \InkAppendixOriginal
  \global\def\InkPageNumberStyle{inkappendixframenumber}}

% -----------------------------------------------------------------------------
% 4. 目录
% -----------------------------------------------------------------------------
% 一级目录带编号（1. 研究背景）；二级目录用紧凑自定义模板：左缩进 2em，
% 形如「3.1 小节名」，避免默认模板占用的额外竖直空间。
\setbeamertemplate{section in toc}[sections numbered]
\setbeamertemplate{subsection in toc}{%
  \leavevmode\leftskip=2em\relax
  \inserttocsectionnumber.\inserttocsubsectionnumber\enspace
  \inserttocsubsection\par}

% -----------------------------------------------------------------------------
% 5. 章节与小节过渡页
% -----------------------------------------------------------------------------
% 写在 frame 之外的 \section / \subsection 会自动插入一张居中过渡页，
% 无需手写额外 frame。两类过渡页都不显示页码（noframenumbering），
% 因此不会占用正文页码，也不受附录编号影响。
% 带星号的 \section* / \subsection* 不触发过渡页。

% 章节页：水平垂直居中，显示「编号 + 标题」，例如「1. 研究背景」
\AtBeginSection[]{%
  \begin{frame}[plain,c,noframenumbering]
    \centering
    {\usebeamerfont{section title}\usebeamercolor[fg]{section title}%
      \insertsectionnumber.\enspace\insertsection\par}
  \end{frame}
}

% 小节页：上下结构。上方灰色「章节号 + 章节名」，下方较大的「章节号.小节号 + 小节名」
\AtBeginSubsection[]{%
  \begin{frame}[plain,c,noframenumbering]
    \centering
    {\normalsize\color{AcademicGray}%
      \insertsectionnumber.\enspace\insertsection\par}
    \vspace{0.8em}
    {\usebeamerfont{section title}\usebeamercolor[fg]{section title}\insertsectionnumber.\insertsubsectionnumber\enspace\insertsubsection\par}
  \end{frame}
}
